\documentclass{article}
\usepackage[T1]{fontenc}
\usepackage{lmodern}
\usepackage{graphicx}
\usepackage{amsmath,amssymb}
\usepackage[a4paper,margin=2cm]{geometry}
\usepackage[absolute,overlay]{textpos}
\setlength{\TPHorizModule}{1mm}
\setlength{\TPVertModule}{1mm}
\usepackage[indonesian]{babel}
\usepackage{hyperref}
\setlength{\emergencystretch}{2em}
% unit: Halaman 23

\begin{document}

% === Halaman 23 (python/native) ===

1. Key-Recovery Attack

• Tujuan: memperoleh kunci kriptografi yang digunakan untuk mengenkripsi/dekripsi.

• Misalkan sebuah sistem menggunakan 

             C = EK(P)  ,     P = plainteks,  K = kunci,  C = cipherteks 

• Dalam key-recovery attack, penyerang berusaha menemukan K.

• Contoh: Misalkan sebuah cipher menggunakan kunci 16 karakter. Penyerang 

melakukan pencarian kunci secara brute force:

            K₁ → dekripsi C → P₁

            K₂ → dekripsi C → P₂

            K₃ → dekripsi C → P₃

               ⋮

   sampai ditemukan kunci yang menghasilkan plainteks yang masuk akal.

23

\end{document}
